% Source: material deferred from the main text.
% A supplies the proof of Theorem 11.7 of Fukuda (2026, p. 419), stated as
% \autoref{thm:implicit-function} in Chapter 8 and there attributed to
% Rudin (1976, Thm 9.28); the construction is the contraction-mapping argument
% advertised on Gerzensee "Real Analysis" slide 65.
% B collects the matrix-differentiation identities used in Chapters 8, 11 and 14,
% several of which appear on Gerzensee "Calculus" slides 45-54 without derivation.
% C gives the multivariate normal distribution and its conditional law, of which
% the bivariate and scalar results of Chapter 14 are special cases; see
% Fukuda (2026), Sec. 9.6.4 (pp. 326-328).
% External references actually cited in this chapter: Ash-Dol{\'e}ans-Dade
% (1999), Axler (2024).

\chapter{The Implicit Function Theorem}\label{app:ift}

\autoref{thm:implicit-function} was stated in \autoref{ch:calculusn} and its
proof deferred. It is given here in full, because the theorem carries the
comparative statics of every model in this volume and because the argument is the
most substantial application of \autoref{thm:banach-fixed-point} outside dynamic
programming.

\section{Two Lemmas}\label{sec:app-ift-lemmas}

\begin{lemma}[Neumann invertibility]\label{lem:neumann}
	Let $M\in M_{m\times m}(\R)$ satisfy $\norm{I-M}<1$ in the operator norm. Then
	$M$ is invertible, $M^{-1}=\sum_{j=0}^{\infty}(I-M)^{j}$, and
	$\norm{M^{-1}}\leq\bigl(1-\norm{I-M}\bigr)^{-1}$.
\end{lemma}

\begin{proof}
	Write $N\eqdef I-M$ and $q\eqdef\norm{N}<1$. Submultiplicativity of the operator
	norm gives $\norm{N^{j}}\leq q^{j}$, so the partial sums
	$S_k\eqdef\sum_{j=0}^{k}N^{j}$ satisfy
	$\norm{S_k-S_l}\leq\sum_{j=l+1}^{k}q^{j}\leq q^{l+1}/(1-q)$ for $k>l$ and are
	therefore Cauchy. The space $M_{m\times m}(\R)$, identified with $\R^{m^{2}}$,
	is complete by \autoref{thm:rn-complete}, so $S_k\to S$ for some matrix $S$
	with $\norm{S}\leq\sum_{j\geq0}q^{j}=(1-q)^{-1}$. Since
	$S_k(I-N)=(I-N)S_k=I-N^{k+1}$ and $\norm{N^{k+1}}\leq q^{k+1}\to0$, passing to
	the limit gives $SM=MS=I$.
\end{proof}

\begin{lemma}[Mean value inequality]\label{lem:mvi}
	Let $U\subseteq\R^{m}$ be open and convex and let $\Phi\colon U\to\R^{m}$ be
	continuously differentiable with $\norm{D\Phi(z)}\leq\kappa$ for every $z\in U$.
	Then $\norm{\Phi(z_1)-\Phi(z_2)}\leq\kappa\norm{z_1-z_2}$ for all
	$z_1,z_2\in U$.
\end{lemma}

\begin{proof}
	Convexity places the segment $s\mapsto z_2+s(z_1-z_2)$, $s\in[0,1]$, inside $U$.
	The map $s\mapsto\Phi\bigl(z_2+s(z_1-z_2)\bigr)$ is continuously differentiable
	with derivative $D\Phi\bigl(z_2+s(z_1-z_2)\bigr)(z_1-z_2)$ by
	\autoref{thm:chain-rule}, so applying \autoref{thm:ftc} to each coordinate,
	\[
		\Phi(z_1)-\Phi(z_2)
		=\int_{0}^{1}D\Phi\bigl(z_2+s(z_1-z_2)\bigr)(z_1-z_2)\,\diff s .
	\]
	Taking norms and using the triangle inequality for integrals gives the bound.
	The scalar mean value theorem (\autoref{thm:mvt}) cannot be applied coordinate
	by coordinate here, since each coordinate would produce its own intermediate
	point and no single point serves for all of them; the integral form avoids the
	difficulty, and that is why it is stated separately.
\end{proof}

\section{Proof of the Theorem}\label{sec:app-ift-proof}

\begin{proof}[Proof of \autoref{thm:implicit-function}]
	Write $A\eqdef\nabla_yf(\bar x,\bar y)$, invertible by hypothesis, and define
	\begin{equation}\label{eq:ift-phi}
		\Phi(x,y)\eqdef y-A^{-1}f(x,y),\qquad(x,y)\in\Lambda .
	\end{equation}
	Because $A^{-1}$ is injective, $f(x,y)=0$ if and only if $\Phi(x,y)=y$: the
	solutions of the system are the fixed points of $\Phi(x,\cdot)$.

	\emph{Step 1: $\Phi(x,\cdot)$ is a contraction.} Differentiating
	\eqref{eq:ift-phi} in $y$ gives
	$\nabla_y\Phi(x,y)=I-A^{-1}\nabla_yf(x,y)$, which vanishes at
	$(\bar x,\bar y)$. Since $f$ is continuously differentiable, $\nabla_yf$ is
	continuous, so there are $r>0$ and $\delta>0$ with
	$\closure{\ball{\delta}{\bar x}}\times\closure{\ball{r}{\bar y}}\subseteq\Lambda$
	and
	\begin{equation}\label{eq:ift-half}
		\norm{I-A^{-1}\nabla_yf(x,y)}\leq\tfrac12
		\qquad\text{whenever }\norm{x-\bar x}\leq\delta,\ \norm{y-\bar y}\leq r .
	\end{equation}
	\autoref{lem:mvi}, applied on the convex set $\ball{r}{\bar y}$ and extended to
	its closure by continuity, gives
	\begin{equation}\label{eq:ift-contraction}
		\norm{\Phi(x,y_1)-\Phi(x,y_2)}\leq\tfrac12\norm{y_1-y_2},
		\qquad y_1,y_2\in\closure{\ball{r}{\bar y}},\ \norm{x-\bar x}\leq\delta .
	\end{equation}

	\emph{Step 2: the ball is invariant.} Since $f(\bar x,\bar y)=0$ we have
	$\Phi(\bar x,\bar y)=\bar y$, and $x\mapsto\Phi(x,\bar y)$ is continuous, so
	after shrinking $\delta$ we may assume
	$\norm{\Phi(x,\bar y)-\bar y}\leq r/2$ for $\norm{x-\bar x}\leq\delta$. For such
	$x$ and any $y\in\closure{\ball{r}{\bar y}}$,
	\[
		\norm{\Phi(x,y)-\bar y}
		\leq\norm{\Phi(x,y)-\Phi(x,\bar y)}+\norm{\Phi(x,\bar y)-\bar y}
		\leq\tfrac12\norm{y-\bar y}+\tfrac{r}{2}\leq r,
	\]
	so $\Phi(x,\cdot)$ maps $\closure{\ball{r}{\bar y}}$ into itself.

	\emph{Step 3: existence and uniqueness of $e$.} The set
	$\closure{\ball{r}{\bar y}}$ is closed in the complete space $\R^{m}$, hence
	complete by \autoref{prop:closed-subset-complete}. By
	\eqref{eq:ift-contraction} and Step 2, $\Phi(x,\cdot)$ is a contraction of
	modulus $\tfrac12$ of that set into itself, so
	\autoref{thm:banach-fixed-point} supplies a unique fixed point, called $e(x)$.
	Setting $\Gamma\eqdef\ball{\delta}{\bar x}$ gives
	$e\colon\Gamma\to\closure{\ball{r}{\bar y}}$ with $f(x,e(x))=0$, which is (i);
	uniqueness of the fixed point at $x=\bar x$ forces $e(\bar x)=\bar y$, which is
	(ii). Uniqueness also yields the final claim of the theorem: inside
	$\Gamma\times\ball{r}{\bar y}$ the solution set of $f=0$ is exactly the graph of
	$e$.

	\emph{Step 4: (iii).} Estimate \eqref{eq:ift-half} says
	$\norm{I-A^{-1}\nabla_yf(x,y)}\leq\tfrac12<1$, so \autoref{lem:neumann} makes
	$A^{-1}\nabla_yf(x,y)$ invertible and hence $\nabla_yf(x,y)$ invertible
	throughout the box, in particular at $(x,e(x))$.

	\emph{Step 5: $e$ is Lipschitz.} For $x_1,x_2\in\Gamma$, using
	$e(x_i)=\Phi(x_i,e(x_i))$ and then \eqref{eq:ift-contraction},
	\[
		\norm{e(x_1)-e(x_2)}
		\leq\underbrace{\norm{\Phi(x_1,e(x_1))-\Phi(x_1,e(x_2))}}
		_{\leq\frac12\norm{e(x_1)-e(x_2)}}
		+\norm{\Phi(x_1,e(x_2))-\Phi(x_2,e(x_2))},
	\]
	and the second term equals
	$\norm{A^{-1}\bigl[f(x_2,e(x_2))-f(x_1,e(x_2))\bigr]}
		\leq\norm{A^{-1}}M\norm{x_1-x_2}$, where $M$ is the supremum of
	$\norm{\nabla_xf}$ over the closed box --- finite by
	\autoref{cor:weierstrass-continuous} --- and \autoref{lem:mvi} is applied in the
	$x$ variable. Rearranging,
	\begin{equation}\label{eq:ift-lipschitz}
		\norm{e(x_1)-e(x_2)}\leq L\norm{x_1-x_2},\qquad L\eqdef2\norm{A^{-1}}M ;
	\end{equation}
	in particular $e$ is continuous.

	\emph{Step 6: differentiability and (iv).} Fix $x\in\Gamma$ and write
	$y\eqdef e(x)$, $B\eqdef\nabla_yf(x,y)$ --- invertible by Step 4 --- and
	$C\eqdef\nabla_xf(x,y)$. For small $h$ put $k\eqdef e(x+h)-e(x)$, so
	$\norm{k}\leq L\norm{h}$ by \eqref{eq:ift-lipschitz}. Differentiability of $f$
	at $(x,y)$ gives
	\[
		0=f(x+h,y+k)-f(x,y)=Ch+Bk+o\bigl(\norm{h}+\norm{k}\bigr),
	\]
	and $\norm{h}+\norm{k}\leq(1+L)\norm{h}$ converts the remainder into
	$o(\norm{h})$. Multiplying by $B^{-1}$,
	\[
		k=-B^{-1}Ch+o(\norm{h}),
	\]
	which is precisely the statement that $e$ is differentiable at $x$ with
	$De(x)=-B^{-1}C$, that is \eqref{eq:ift-derivative}. Continuity of $De$ follows
	because $x\mapsto\nabla_yf(x,e(x))$ and $x\mapsto\nabla_xf(x,e(x))$ are
	continuous, by continuity of $Df$ and of $e$, and because matrix inversion is
	continuous on the invertible matrices --- itself a consequence of
	\autoref{lem:neumann}. Hence $e$ is continuously differentiable.
\end{proof}

\begin{remark}[Deriving the inverse function theorem]\label{rem:ift-implies-inverse}
	\autoref{thm:inverse-function} follows by applying
	\autoref{thm:implicit-function} to $F(z,x)\eqdef g(x)-z$ on
	$\R^{n}\times A$ at the point $\bigl(g(\bar x),\bar x\bigr)$: here the
	derivative in the second block is $\nabla_xF=Dg(\bar x)$, nonsingular by
	hypothesis, so there is a continuously differentiable $e$ with $g(e(z))=z$ near
	$g(\bar x)$, and \eqref{eq:ift-derivative} gives
	$De(z)=-\bigl[Dg(e(z))\bigr]^{-1}(-I)=\bigl[Dg(e(z))\bigr]^{-1}$.
\end{remark}

\chapter{Matrix Differentiation}\label{app:matrix-calculus}

The identities below are used repeatedly in
Chapters~\ref{ch:calculusn},~\ref{ch:projection} and~\ref{ch:probability}. All
vectors are columns, gradients of scalar functions are columns, and $Df$ denotes
the Jacobian, so that $Df(x)=\nabla f(x)^{\transp}$ for scalar $f$. This
convention is the one fixed in the front matter, and it is the only source of
transposition errors in what follows.

\section{Linear and Quadratic Forms}\label{sec:app-quadratic}

\begin{proposition}\label{prop:matrix-derivatives}
	Let $a\in\R^{n}$, $A\in M_{n\times n}(\R)$, $B\in M_{m\times n}(\R)$,
	$x\in\R^{n}$ and $y\in\R^{n}$.
	\begin{enumerate}[(i)]
		\item $\nabla\bigl(\ip{a}{x}\bigr)=a$ and $D(Bx)=B$.
		\item $\nabla\bigl(x^{\transp}Ax\bigr)=(A+A^{\transp})x$; in particular
		      $\nabla\bigl(x^{\transp}Ax\bigr)=2Ax$ when $A$ is symmetric, and
		      $\nabla\norm{x}^{2}=2x$.
		\item $\Hess\bigl(x^{\transp}Ax\bigr)=A+A^{\transp}$, a constant matrix.
		\item If $\varphi(x)\eqdef f(Bx)$ with $f$ differentiable, then
		      $\nabla\varphi(x)=B^{\transp}\nabla f(Bx)$.
		\item $\nabla_x\bigl(x^{\transp}Ay\bigr)=Ay$ and
		      $\nabla_y\bigl(x^{\transp}Ay\bigr)=A^{\transp}x$.
	\end{enumerate}
\end{proposition}

\begin{proof}
	(i) $\ip{a}{x}=\sum_ia_ix_i$ has $\partial/\partial x_j=a_j$; the $i$-th
	coordinate of $Bx$ is $\sum_jB_{ij}x_j$, whose partial derivative in $x_j$ is
	$B_{ij}$.

	(ii) Writing $q(x)=\sum_{i,j}A_{ij}x_ix_j$,
	\[
		\pderiv{q}{x_k}=\sum_{j}A_{kj}x_j+\sum_{i}A_{ik}x_i
		=\bigl(Ax\bigr)_k+\bigl(A^{\transp}x\bigr)_k,
	\]
	the two sums arising from the occurrences of $x_k$ in the first and in the
	second factor. Taking $A=I$ gives $\nabla\norm{x}^{2}=2x$.

	(iii) Differentiate (ii) once more: $x\mapsto(A+A^{\transp})x$ is linear with
	Jacobian $A+A^{\transp}$ by (i).

	(iv) \autoref{thm:chain-rule} gives $D\varphi(x)=Df(Bx)B$, and transposing,
	$\nabla\varphi(x)=B^{\transp}\nabla f(Bx)$. The transpose is where the
	convention matters: $Df$ is a row and $\nabla f$ a column.

	(v) $x^{\transp}Ay=\sum_{i,j}x_iA_{ij}y_j$ is linear in $x$ for fixed $y$ with
	coefficient vector $Ay$, and linear in $y$ for fixed $x$ with coefficient
	vector $A^{\transp}x$; apply (i).
\end{proof}

\begin{eg}[Where each identity is used]\label{eg:matrix-uses}
	Identity (ii) with symmetric $A$ produced \eqref{eq:quadratic-gradient} and the
	second-order conditions of \autoref{thm:second-order-conditions}. Identity (iv)
	with $f=\norm{\cdot}^{2}$ gives the least-squares normal equations: minimising
	$\norm{y-X\beta}^{2}$ over $\beta$ has gradient
	$-2X^{\transp}(y-X\beta)$, whose vanishing is
	$X^{\transp}X\beta=X^{\transp}y$, the calculus route to
	\autoref{prop:ols-calculus} and the algebraic counterpart of the projection
	argument of \autoref{prop:ols-projection}. Identity (iii) shows that the Hessian
	of a quadratic form is constant, so a quadratic objective is concave globally or
	not at all --- which is why \autoref{thm:hessian-convexity} needs to be checked
	only once for such objectives rather than at every point.
\end{eg}

\section{Inverse, Trace and Determinant}\label{sec:app-inverse}

\begin{proposition}\label{prop:matrix-derivatives-2}
	Let $X\in M_{n\times n}(\R)$ be invertible and $A\in M_{n\times n}(\R)$.
	\begin{enumerate}[(i)]
		\item If $t\mapsto X(t)$ is differentiable with $X(t)$ invertible, then
		      \begin{equation}\label{eq:deriv-inverse-matrix}
			      \deriv{}{t}X(t)^{-1}
			      =-X(t)^{-1}\Bigl(\deriv{}{t}X(t)\Bigr)X(t)^{-1}.
		      \end{equation}
		\item $\partial\tr(AX)/\partial X=A^{\transp}$, entrywise.
		\item $\partial\log\det X/\partial X=(X^{-1})^{\transp}$.
	\end{enumerate}
\end{proposition}

\begin{proof}
	(i) Differentiate $X(t)X(t)^{-1}=I$ by the product rule
	(\autoref{prop:deriv-product}, applied entrywise), obtaining
	$\dot XX^{-1}+X\deriv{}{t}(X^{-1})=0$, and multiply on the left by $X^{-1}$.

	(ii) $\tr(AX)=\sum_{i,j}A_{ij}X_{ji}$, whose partial derivative in $X_{ji}$ is
	$A_{ij}$; assembling these into the matrix indexed by $(j,i)$ gives
	$A^{\transp}$.

	(iii) Jacobi's formula $\diff(\det X)=\det X\cdot\tr\bigl(X^{-1}\,\diff X\bigr)$,
	a consequence of the cofactor expansion of \autoref{def:determinant}, gives
	$\diff\log\det X=\tr\bigl(X^{-1}\,\diff X\bigr)$; comparing with (ii) at
	$A=X^{-1}$ gives the stated gradient. A derivation of Jacobi's formula is in
	\textcite[Ch.~10]{axler2024linear}.
\end{proof}

Identity \eqref{eq:deriv-inverse-matrix} is the matrix analogue of
$\deriv{}{t}(1/x)=-\dot x/x^{2}$, and it is what the Riccati recursion
\eqref{eq:riccati-recursion} and the comparative statics formula
\eqref{eq:ift-derivative} reduce to in the scalar case. Identity (iii) is the
first-order condition of Gaussian maximum likelihood with respect to a covariance
matrix, and it is the reason the log-likelihood is concave in the precision
matrix rather than in the covariance matrix.

\chapter{The Multivariate Normal Distribution}\label{app:normal}

\autoref{thm:normal-updating} and \autoref{prop:bivariate-normal} are the scalar
and bivariate cases of one statement, recorded here in general form because
econometric and macroeconomic applications use it with several signals and
several states at once.

\section{Definition and Affine Images}\label{sec:app-mvn}

\begin{definition}[Multivariate normal]\label{def:mvn}
	A random vector $X=(X_1,\dots,X_n)^{\transp}$ is \dfnas{multivariate
		normal}{multivariate normal distribution} with mean $\mu\in\R^{n}$ and
	covariance $\Sigma\succ0$, written $X\sim\mathcal{N}(\mu,\Sigma)$, if it has
	density
	\begin{equation}\label{eq:mvn-density}
		f(x)=\frac{1}{(2\pi)^{n/2}(\det\Sigma)^{1/2}}
		\exp\Bigl(-\tfrac12(x-\mu)^{\transp}\Sigma^{-1}(x-\mu)\Bigr),
		\qquad x\in\R^{n}.
	\end{equation}
	\nidx{mvn}{$\mathcal{N}(\mu,\Sigma)$, multivariate normal distribution}
\end{definition}

That \eqref{eq:mvn-density} integrates to one is
\autoref{prop:gaussian-integral} combined with the substitution
$z=\Sigma^{-1/2}(x-\mu)$, whose Jacobian determinant is
$(\det\Sigma)^{-1/2}$ by \autoref{thm:change-of-variables}; the symmetric square
root $\Sigma^{1/2}=Q\Lambda^{1/2}Q^{\transp}$ exists by \autoref{thm:spectral}
and has strictly positive diagonal $\Lambda^{1/2}$ by
\autoref{prop:definiteness-eigen}.

\begin{proposition}[Affine images and the moment generating function]\label{prop:mvn-affine}
	Let $X\sim\mathcal{N}(\mu,\Sigma)$, $b\in\R^{m}$ and $B\in M_{m\times n}(\R)$ of
	rank $m$. Then
	\begin{equation}\label{eq:mvn-affine}
		b+BX\sim\mathcal{N}\bigl(b+B\mu,\ B\Sigma B^{\transp}\bigr),
		\qquad
		\E\bigl[e^{\ip{t}{X}}\bigr]
		=\exp\Bigl(\ip{t}{\mu}+\tfrac12t^{\transp}\Sigma t\Bigr)
		\quad\text{for }t\in\R^{n}.
	\end{equation}
	In particular every linear combination $\ip{a}{X}$ is univariate normal with
	mean $\ip{a}{\mu}$ and variance $a^{\transp}\Sigma a$.
\end{proposition}

\begin{proof}
	For the moment generating function, complete the square in the exponent of
	\eqref{eq:mvn-density}: with $\nu\eqdef\mu+\Sigma t$,
	\[
		\ip{t}{x}-\tfrac12(x-\mu)^{\transp}\Sigma^{-1}(x-\mu)
		=\ip{t}{\mu}+\tfrac12t^{\transp}\Sigma t
		-\tfrac12(x-\nu)^{\transp}\Sigma^{-1}(x-\nu),
	\]
	as multiplying out both sides and using the symmetry of $\Sigma^{-1}$ verifies.
	Integrating the remaining Gaussian kernel gives $1$ and leaves the stated
	formula. For the affine image, the moment generating function of $b+BX$ at
	$s\in\R^{m}$ is
	\[
		\E\bigl[e^{\ip{s}{b+BX}}\bigr]
		=e^{\ip{s}{b}}\,\E\bigl[e^{\ip{B^{\transp}s}{X}}\bigr]
		=\exp\Bigl(\ip{s}{b+B\mu}+\tfrac12s^{\transp}B\Sigma B^{\transp}s\Bigr),
	\]
	which is the moment generating function of
	$\mathcal{N}(b+B\mu,B\Sigma B^{\transp})$; a distribution on $\R^{m}$ whose
	moment generating function is finite on a neighbourhood of the origin is
	determined by it \autocite[Ch.~5]{ash1999probability}. The rank condition keeps
	$B\Sigma B^{\transp}$ positive definite.
\end{proof}

\begin{corollary}[Uncorrelated implies independent]\label{cor:mvn-independence}
	Let $(X_1^{\transp},X_2^{\transp})^{\transp}\sim\mathcal{N}(\mu,\Sigma)$, with
	$\Sigma$ partitioned into blocks $\Sigma_{11},\Sigma_{12},\Sigma_{21},
		\Sigma_{22}$ conformably. Then $X_1$ and $X_2$ are independent if and only if
	$\Sigma_{12}=0$.
\end{corollary}

\begin{proof}
	Independence implies zero covariance for any square-integrable pair, by
	\autoref{prop:expectation-properties}(iv). Conversely, if $\Sigma_{12}=0$ then
	$\Sigma$ and $\Sigma^{-1}$ are block diagonal and
	$\det\Sigma=\det\Sigma_{11}\det\Sigma_{22}$, so \eqref{eq:mvn-density} factors
	into the product of the two marginal densities.
\end{proof}

The corollary is false without joint normality. Let $Z$ be standard normal and
let $S$ be independent of $Z$ with $\Prob(S=1)=\Prob(S=-1)=\tfrac12$. Then $SZ$
is standard normal and
$\Cov(Z,SZ)=\E[S]\E[Z^{2}]=0$, yet $\abs{Z}=\abs{SZ}$, so the two are far from
independent. The pair $(Z,SZ)$ is not jointly normal, since $Z+SZ$ equals $0$
with probability $\tfrac12$ and no non-degenerate normal variable has an atom.

\section{The Conditional Law}\label{sec:app-mvn-conditional}

\begin{theorem}[Conditional distribution]\label{thm:mvn-conditional}
	Let $(X_1^{\transp},X_2^{\transp})^{\transp}$ be multivariate normal with mean
	$(\mu_1^{\transp},\mu_2^{\transp})^{\transp}$ and covariance blocks
	$\Sigma_{11},\Sigma_{12},\Sigma_{21},\Sigma_{22}$, with $\Sigma_{22}$
	invertible. Then
	\begin{equation}\label{eq:mvn-conditional}
		X_1\mid X_2=x_2\ \sim\
		\mathcal{N}\Bigl(\mu_1+\Sigma_{12}\Sigma_{22}^{-1}(x_2-\mu_2),\
		\Sigma_{11}-\Sigma_{12}\Sigma_{22}^{-1}\Sigma_{21}\Bigr).
	\end{equation}
	The conditional mean is affine in $x_2$ and the conditional covariance does not
	depend on $x_2$.
\end{theorem}

\begin{proof}
	Set $Z\eqdef X_1-\Sigma_{12}\Sigma_{22}^{-1}X_2$. By \autoref{prop:mvn-affine}
	the pair $(Z^{\transp},X_2^{\transp})^{\transp}$ is an affine image of
	$(X_1^{\transp},X_2^{\transp})^{\transp}$ and is therefore jointly normal, and
	\[
		\Cov(Z,X_2)=\Sigma_{12}-\Sigma_{12}\Sigma_{22}^{-1}\Sigma_{22}=0 .
	\]
	\autoref{cor:mvn-independence} makes $Z$ independent of $X_2$. Hence,
	conditionally on $X_2=x_2$, the vector
	$X_1=Z+\Sigma_{12}\Sigma_{22}^{-1}x_2$ is $Z$ shifted by a constant, so its
	conditional law is normal with mean
	\[
		\E[Z]+\Sigma_{12}\Sigma_{22}^{-1}x_2
		=\mu_1-\Sigma_{12}\Sigma_{22}^{-1}\mu_2+\Sigma_{12}\Sigma_{22}^{-1}x_2
	\]
	and covariance
	\[
		\Var(Z)=\Sigma_{11}-\Sigma_{12}\Sigma_{22}^{-1}\Sigma_{21}
		-\Sigma_{12}\Sigma_{22}^{-1}\Sigma_{21}
		+\Sigma_{12}\Sigma_{22}^{-1}\Sigma_{22}\Sigma_{22}^{-1}\Sigma_{21}
		=\Sigma_{11}-\Sigma_{12}\Sigma_{22}^{-1}\Sigma_{21},
	\]
	which together give \eqref{eq:mvn-conditional}.
\end{proof}

\begin{remark}[Recovering the results of \autoref{ch:probability}]
	\label{rem:mvn-specialisation}
	Take the blocks to be scalars with $X_1=Y$ and $X_2=X$. Then
	$\Sigma_{12}\Sigma_{22}^{-1}=\Cov(X,Y)/\Var(X)$ and
	\eqref{eq:mvn-conditional} is \eqref{eq:gaussian-conditional}, so
	\autoref{prop:bivariate-normal} is the bivariate case.

	Take instead $X_1=X\sim\mathcal{N}(\mu_x,\sigma_x^{2})$ and
	$X_2=Y=X+\varepsilon$ with
	$\varepsilon\sim\mathcal{N}(0,\sigma_\varepsilon^{2})$ independent of $X$. The
	pair is jointly normal by \autoref{prop:mvn-affine}, with
	$\Sigma_{11}=\Sigma_{12}=\sigma_x^{2}$ and
	$\Sigma_{22}=\sigma_x^{2}+\sigma_\varepsilon^{2}$, so
	\eqref{eq:mvn-conditional} gives conditional mean
	$\mu_x+\frac{\sigma_x^{2}}{\sigma_x^{2}+\sigma_\varepsilon^{2}}(y-\mu_x)$ and
	conditional variance
	$\sigma_x^{2}-\frac{\sigma_x^{4}}{\sigma_x^{2}+\sigma_\varepsilon^{2}}
		=\frac{\sigma_x^{2}\sigma_\varepsilon^{2}}
		{\sigma_x^{2}+\sigma_\varepsilon^{2}}$, which is
	\eqref{eq:normal-updating}. The completing-the-square derivation given there is
	therefore an instance of the orthogonalisation argument used above, and the
	statement that precisions add is the identity
	$\bigl(\Sigma_{11}-\Sigma_{12}\Sigma_{22}^{-1}\Sigma_{21}\bigr)^{-1}
		=\sigma_x^{-2}+\sigma_\varepsilon^{-2}$ in this parametrisation.
\end{remark}
